\documentclass[a4paper, 10pt]{article}
\usepackage[margin=1in]{geometry}
\usepackage{amsfonts, amsmath, amssymb, amsthm}
%\usepackage[none]{hyphenat}
\usepackage[T2A]{fontenc}
\usepackage{fancyhdr} %create a custom header and footer
\usepackage[utf8]{inputenc}
\usepackage[english, main=ukrainian]{babel}
\usepackage{pgfplots}
\pgfplotsset{compat = newest}
\usepackage{changepage}

\usepackage{nccmath}
\usepackage{gnuplottex}


\usepackage{enumitem}
\usepgfplotslibrary{fillbetween}
\usepackage{tikz}
\usepackage{graphicx}
\usepackage{caption}
\usepackage{float}
\usepackage{physics}
\usepackage[unicode]{hyperref}
\usepackage{pdfpages}
\usepackage{tikz-3dplot}
\usepackage{bbm}
\usepackage{tikz-cd}

\usetikzlibrary{spy,angles,quotes}

\fancyhead{}
\fancyfoot{}
\parindent 0ex
\DeclareMathOperator*\uplim{\overline{lim}}
\DeclareMathOperator*\downlim{\underline{lim}}
\DeclareMathOperator{\wordgrad}{grad}
\def\stackbelow#1#2{\underset{\displaystyle\overset{\displaystyle\shortparallel}{#2}}{#1}}
\def\departial#1#2{\dfrac{\partial {#1}}{\partial {#2}}}
\def\seconddepartial#1#2#3{\ifthenelse{\equal{#2}{#3}}{\dfrac{\partial^2 {#1}}{\partial {#2}^2}}{\dfrac{\partial^2 {#1}}{\partial {#2} \partial {#3}}}}
\def\huge{\displaystyle}
\def\bigline{\vspace{5mm}\\}

\def\qed{$\blacksquare$}

\def\rightproof{$\boxed{\Rightarrow}$ }
\def\leftproof{$\boxed{\Leftarrow}$ }


\def\noProof{\\ \textit{Без доведення.}}
\DeclareMathOperator\sign{sgn}

\newtheoremstyle{theoremdd}% name of the style to be used
  {\topsep}% measure of space to leave above the theorem. E.g.: 3pt
  {\topsep}% measure of space to leave below the theorem. E.g.: 3pt
  {\normalfont}% name of font to use in the body of the theorem
  {0pt}% measure of space to indent
  {\bfseries}% name of head font
  {}% punctuation between head and body
  { }% space after theorem head; " " = normal interword space
  {\thmname{#1}\thmnumber{ #2}\textnormal{\thmnote{ \textbf{#3}\\}}}

\theoremstyle{theoremdd}
\newtheorem{theorem}{Theorem}[subsection]

\theoremstyle{theoremdd}
\newtheorem*{theorem*}{Theorem.}
 
\theoremstyle{theoremdd}
\newtheorem{definition}[theorem]{Definition}

\theoremstyle{theoremdd}
\newtheorem{samedef}[theorem]{Definition}

\theoremstyle{theoremdd}
\newtheorem{example}[theorem]{Example}

\theoremstyle{theoremdd}
\newtheorem{proposition}[theorem]{Proposition}

\theoremstyle{theoremdd}
\newtheorem{remark}[theorem]{Remark}

\theoremstyle{theoremdd}
\newtheorem{lemma}[theorem]{Lemma}

\theoremstyle{theoremdd}
\newtheorem{corollary}[theorem]{Corollary}

\newcommand\thref[1]{\textbf{Th.~\ref{#1}}}
\newcommand\defref[1]{\textbf{Def.~\ref{#1}}}
\newcommand\exref[1]{\textbf{Ex.~\ref{#1}}}
\newcommand\prpref[1]{\textbf{Prp.~\ref{#1}}}
\newcommand\rmref[1]{\textbf{Rm.~\ref{#1}}}
\newcommand\lmref[1]{\textbf{Lm.~\ref{#1}}}
\newcommand\crlref[1]{\textbf{Crl.~\ref{#1}}}

\makeatletter
\renewenvironment{proof}[1][Proof.\\]{\par
\pushQED{\hfill \qed}%
\normalfont \topsep6\p@\@plus6\p@\relax
\trivlist
\item\relax
{\bfseries
#1\@addpunct{.}}\hspace\labelsep\ignorespaces
}{%
\popQED\endtrivlist\@endpefalse
}
\makeatother

\newenvironment{pfMI}{\vspace*{-3mm} \textbf{\\ Proof MI. \\}}{\hfill $\blacksquare$}
\newenvironment{pfNoTh}{\textbf{Proof. \\}}{$\blacksquare$}

\newcommand\Norm[1]{\|#1\|}
\newcommand{\notimplies}{\;\not\!\!\!\implies}
\newcommand{\tounif}{^\rightarrow_\rightarrow}

\DeclareMathOperator{\Mat}{Mat}
\DeclareMathOperator{\Cl}{Cl}
\DeclareMathOperator{\Int}{Int}
\DeclareMathOperator{\sgn}{sgn}
\DeclareMathOperator{\linspan}{span}
%\DeclareMathOperator{\gradient}{grad}

\begin{document}
\section*{Вступ до диференціальних форм. $0$-форма}
Мета теорії диференціальних форм: узагальнення диференціальне та інтегральне числення на 'безкоординатний' випадок (тобто на многовидах).\\
Для спрощенного розуміння поки розглянемо випадок на $\mathbb{R}^3$, бо ми щось про диференціальне та інтегральне числення там знаємо. Координати будемо позначати $x_1,x_2,x_3$. Поступово все будемо узагальнювати.
\bigskip \\
Нехай $X \subset \mathbb{R}^3$ -- деяка відкрита підмножина.\\
\textbf{$0$-формою} будемо називати ось такий простір: $\Omega^0(X) \overset{\text{def.}}{=} C^\infty(X)$. По суті, це просто набір гладких функцій трьох змінних. Нічого особливого.
\bigskip \\
Мабуть, перед цим треба буде трошки поговорити за новий апарат лінійної алгеби.
\newpage

\section{Зовнішня алгебра}
Нехай $V$ -- векторний простір та $U \subset V$ -- двовимірний підпростір.
\bigskip \\
Треба пригадати, що таке \textbf{кососиметрична} форма (її англомовним чином називають \textbf{alternating} map). Простір цих кососиметричних форм позначають за $A_k(V)$.

\begin{definition}
\textbf{Другим зовнішнім степенем} скінченновимірного векторного простору називають дуальний простір $(A_2(V))^*$.\\
Позначення: $\Lambda^2 V$.\\
Елементами $\Lambda^2 V$ називають \textbf{$2$-векторами}.
\end{definition}

\begin{definition}
Нехай $u,v \in V$.\\
\textbf{Зовнішнім добутком} $u \wedge v \in \Lambda^2 V$ називають наступне:
\begin{align*}
(u \wedge v)(B) = B(u,v),\ B \in A_2(V)
\end{align*}
\end{definition}

\begin{proposition}[Властивості зовнішнього добутку]
Справедливе наступне:
\begin{enumerate}[nosep,wide=0pt,label={\arabic*)}]
\item $v \wedge u = -u \wedge v$;
\item $(\lambda_1 u_1 + \lambda_2 u_2) \wedge v = \lambda u_1 \wedge v + \lambda u_2 \wedge u$;
\item Якщо $\{v_1,\dots,v_n\}$ -- базис $V$, тоді $\{ v_i \wedge v_j \mid i < j \}$ -- базис $\Lambda^2 V$.
\end{enumerate}
\end{proposition}

\begin{proposition}
Нехай $u \in V$ -- ненулевий вектор. Тоді $u \wedge v = 0 \iff v = \lambda u$ для деякого скаляра $\lambda$.
\end{proposition}

\newpage
\section{$1$-форми}
\subsection{Означення $1$-фоми}
Ми визначимо деякі об'єкти: $dx_1,\ dx_2,\ dx_3 \colon X \to \mathbb{R}$. Кожний з них представляє наступне:
\begin{align*}
dx_i(\vec{x}) = x_i
\end{align*}
Чому саме так? Well, якщо глянути на функцію $f(x_1,x_2,x_3) = x_1$, то коли будемо рахувати диференціал, отримуємо $df(x_1,x_2,x_3) = $
Разом система $\{dx_1,\ dx_2,\ dx_3\}$ формує базис на $X^*$. Чому такі позначення (тобто чому вони збігаються з диференціалами), буде виправдано пізніше.\\
\textbf{$1$-формою} будемо називати ось такий простір:
\begin{align*}
\Omega^1(X) \overset{\text{def.}}{=} \{ f_1\,dx_1 + f_2\,dx_2 + f_3\,dx_3 \mid f_1,f_2,f_3 \in C^\infty(X) \}
\end{align*}
Коротко записуючи, $\Omega^1(X) = \linspan_{C^\infty(X)}\{dx_1,\, dx_2,\, dx_3\}$ -- лінійна оболонка таких елементів, де коефіцієнтами будуть гладкі функції. Подібний об'єкт зустрічали, коли обчислювали так званий \textit{диференціал функції $f$, що визначали як $df = f'_x\,dx + f'_y\,dy + f'_z\,dz$.}
\bigskip \\
Природно ми можемо визначити відображення $d \colon \Omega^0(X) \to \Omega^1(X)$ ось таким чином: $f \mapsto df = f'_{x_1}\,dx_1 + f'_{x_2}\,dx_2 + f'_{x_3}\,dx_3$. Дане відображення буде природно назвати \textbf{диференціалом}.

\subsection{$2$-форми}
Ми визначимо деякі інші вже об'єкти: $dx_1 \wedge dx_2,\ dx_1 \wedge dx_3,\ dx_2 \wedge dx_3$. Що це, щас буде.\\
По суті, ми із $\Omega^1(X)$ забираємо всі базисні вектори $dx_1\,dx_2\,dx_3$ та проводимо попарно операцію $\wedge$ (англійською ця хрінь називається \textbf{wedge}). Така операція задовлоьнятиме властивості:
\begin{align*}
dx_2 \wedge dx_1 = -dx_1 \wedge dx_2 \text{ (із іншими аналогічно) } \\
dx_1 \wedge dx_1 = 0 \text{ (із іншими аналогічно) }
\end{align*}
Тобто дана операція над собою онулює, а також можна змінювати порядок зі знаком мінус.\\
Тепер побудуємо векторний простір, де $\{dx_1 \wedge dx_2,\ dx_1 \wedge dx_3,\ dx_2 \wedge dx_3\}$ формує базис. Ми будемо записувати $dx_i$ в порядку зростання для зручності.\\
\textbf{$2$-формою} будемо називати ось такий простір:\\
$\Omega^2(X) \overset{\text{def.}}{=} \{ f_3\, dx_1 \wedge dx_2 + f_2\, dx_1 \wedge dx_3 + f_1\, dx_2 \wedge dx_3 \mid f_1,f_2,f_3 \in C^\infty(X) \}$.
\bigskip \\
Хочемо аналогічно визначити відображення $d \colon \Omega^1(X) \to \Omega^0(X)$.
\end{document}